% Part:sets-functions-relations
% Chapter: size-of-sets
% Section: reduction-alt

\documentclass[../../../include/open-logic-section]{subfiles}

\begin{document}

\olfileid{sfr}{siz}{red-alt}

\olsection{হ্রাসকরণ}

\begin{editorial}
  এই অংশে হ্রাসকরণের সাহায্যে অতালিকায়নযোগ্যতা প্রমাণ করা হয়েছে; ফলগুলি \olref[nen-alt]{sec}-এর ফলের সঙ্গে মেলে। \olref[nen]{sec}-এর ফলের সঙ্গে মেলে এমন একটি বিকল্প, কিছুটা বিস্তৃত সংস্করণ \olref[red]{sec}-এ দেওয়া আছে।
\end{editorial}

কর্ণীকরণ যুক্তির সাহায্যে আমরা প্রমাণ করেছি যে $\Bin^\omega$ !!{nonenumerable}। একই ধরনের কর্ণীকরণ যুক্তি দিয়ে দেখিয়েছি যে $\Pow{\Nat}$ !!{nonenumerable}। কিন্তু $\Pow{\Nat}$ যে !!{nonenumerable}, তা প্রমাণের আরেকটি উপায় এখানে: দেখাও যে \emph{$\Pow{\Nat}$ যদি !!{enumerable} হয়, তবে $\Bin^\omega$-ও !!{enumerable}}। যেহেতু আমরা জানি $\Bin^\omega$ !!{nonenumerable}, তাই এর থেকে অনুসৃত হবে যে $\Pow{\Nat}$-ও তাই।

একে একটি সমস্যাকে অন্য সমস্যায় \emph{হ্রাস করা} বলা হয়। এই ক্ষেত্রে আমরা $\Bin^\omega$-কে তালিকায়ন করার সমস্যাটিকে $\Pow{\Nat}$-কে তালিকায়ন করার সমস্যায় হ্রাস করি। পরেরটির একটি সমাধান—$\Pow{\Nat}$-এর একটি তালিকায়ন—আগেরটির একটি সমাধান—$\Bin^\omega$-এর একটি তালিকায়ন—দিত।

কোনো সেট~$B$-কে তালিকায়ন করার সমস্যাকে সেট~$A$-কে তালিকায়ন করার সমস্যায় হ্রাস করতে আমরা~$A$-এর একটি তালিকায়নকে~$B$-এর তালিকায়নে বদলানোর একটি উপায় দিই। এর সবচেয়ে সহজ উপায় হলো !!a{surjection} $f\colon A \to B$ সংজ্ঞায়িত করা। $x_1$, $x_2$, \dots{} যদি~$A$-কে তালিকায়িত করে, তবে $f(x_1)$, $f(x_2)$, \dots{}~$B$-কে তালিকায়িত করবে। আমাদের ক্ষেত্রে আমরা !!a{surjection} $f\colon \Pow{\Nat}
\to \Bin^\omega$ খুঁজছি।

\begin{prob}
দেখাও যে কোনো !!{injective} অপেক্ষক $g\colon B \to A$ থাকলে এবং~$B$ !!{nonenumerable} হলে~$A$-ও তাই।~$A$-এর একটি তালিকায়নকে~$B$-এর তালিকায়নে বদলাতে কীভাবে~$g$ ব্যবহার করা যায় তা দেখিয়ে প্রমাণটি করো।
\end{prob}

\begin{proof}[হ্রাসকরণের সাহায্যে {\olref[nen-alt]{thm:nonenum-pownat}}-এর প্রমাণ]
হ্রাসকরণের জন্য ধরি, $\Pow{\Nat}$ !!{enumerable}, এবং তাই তার একটি তালিকায়ন $N_{1}$, $N_{2}$, $N_{3}$, \dots{} আছে।

$f(N)$-কে সেই প্রতীকক্রম $s$ নিই, যার জন্য $s(n) = 1$ যদি এবং কেবল যদি $n \in N$, এবং অন্যথায় $s(n) = 0$; এভাবে অপেক্ষক $f \colon \Pow{\Nat} \to \Bin^\omega$ সংজ্ঞায়িত করি। % BN-SRC-005 TARGET-CORRECTION-BEGIN
\par\smallskip\noindent\textnormal{\small\textbf{উৎস-সংশোধন (BN-SRC-005 / OLSIZ-010):} হিমায়িত ইংরেজি উৎসে এখানে $s_k$ লেখা হলেও $k$ বাঁধা বা নির্ধারিত নয়। বাংলা পাঠে ফল অনুক্রমটির নাম সর্বত্র $s$ করা হয়েছে; তার $n$-তম মান $n\in N$ হলে $1$, অন্যথায় $0$।} % BN-SRC-005 TARGET-CORRECTION-END

এটি স্পষ্টতই একটি অপেক্ষক সংজ্ঞায়িত করে, কারণ $N \subseteq \Nat$ হলে যেকোনো $n \in \Nat$ হয় $N$-এর !!a{element}, নয়তো নয়। যেমন, জোড় স্বাভাবিক সংখ্যার সেট $2\Nat = \Setabs{2n}{n \in \Nat} = \{0,2, 4, 6,
\dots\}$-কে প্রতীকক্রম $1010101\dots$-এ পাঠানো হয়; $\emptyset$-কে $0000\dots$-এ; এবং $\Nat$-কে
$1111\dots$-এ পাঠানো হয়।

এটি !!{surjective}-ও: $0$ ও $1$ দিয়ে গঠিত প্রত্যেক প্রতীকক্রম স্বাভাবিক সংখ্যার কোনো সেটের সঙ্গে সঙ্গতিপূর্ণ, যথা প্রতীকক্রমটির যেসব স্থানে~$1$ আছে সেই স্থানগুলির সঙ্গে সংশ্লিষ্ট স্বাভাবিক সংখ্যাগুলি যে সেটের সদস্য। আরও নির্দিষ্টভাবে, $s \in \Bin^\omega$ হলে নিচের নিয়মে $N
\subseteq \Nat$ সংজ্ঞায়িত করি:
\[
N = \Setabs{n \in \Nat}{s(n) = 1}
\]
তাহলে~$f$-এর সংজ্ঞা দেখে যাচাই করা যায় যে $f(N) = s$।

এখন তালিকাটি বিবেচনা করো:
\[
f(N_1), f(N_2), f(N_3), \dots
\]
যেহেতু $f$ !!{surjective}, তাই $\Bin^\omega$-এর প্রত্যেক সদস্যকে কোনো আর্গুমেন্টের জন্য~$f$-এর মান হিসেবে আসতে হবে; ফলে সেটি তালিকাতেও আসবে। তাই এই তালিকাকে~$\Bin^\omega$-এর সবকিছু তালিকায়িত করতেই হবে।

সুতরাং $\Pow{\Nat}$ !!{enumerable} হলে $\Bin^\omega$ !!{enumerable} হতো। কিন্তু $\Bin^\omega$ !!{nonenumerable}
(\olref[nen-alt]{thm:nonenum-bin-omega})। তাই $\Pow{\Nat}$ !!{nonenumerable}।
\end{proof}

%\begin{explain}
%হ্রাস কোন অভিমুখে যায়, তা নিয়ে সহজেই বিভ্রান্তি হতে পারে। যেমন, !!a{surjective} অপেক্ষক $g \colon \Bin^\omega \to X$, $X$ যে !!{nonenumerable} তা \emph{প্রতিষ্ঠা করে না}। ($g(s) = s(1)$ দ্বারা নির্ধারিত $g
%\colon \Bin^\omega \to \Bin$ বিবেচনা করো; অপেক্ষকটি $0$ ও $1$-এর অনুক্রমকে তার প্রথম !!{element}-এ পাঠায়। এটি সমাপতিত, কারণ কিছু অনুক্রম $0$ দিয়ে এবং কিছু অনুক্রম $1$ দিয়ে শুরু হয়। কিন্তু $\Bin$ সসীম।) আরও লক্ষ করো, অপেক্ষক~$f$-কে সমাপতিত হতেই হবে, নইলে যুক্তিটি কাজ করে না:
%$f(x_1)$, $f(x_2)$, \dots{}-এ~$Y$-এর সব !!{element}s অবশ্যই থাকবে—এ নিশ্চয়তা তখন পাওয়া যেত না। যেমন, $h\colon \Nat \to
%\Bin^\omega$-কে নিচের মতো সংজ্ঞায়িত করি:
%\[
%h(n) = \underbrace{000\dots0}_{\text{$n$টি $0$}}
%\]
%এটি একটি অপেক্ষক, কিন্তু $\Nat$ !!{enumerable}।
%\end{explain}

\begin{prob}\label{sfr:siz:red:prob:nat-nat}
  হ্রাসকরণ যুক্তির সাহায্যে দেখাও যে সকল অপেক্ষক $f\colon \Nat \to \Nat$-এর সেট~$X$
  !!{nonenumerable}। (ইঙ্গিত: $X$ থেকে~$\Bin^\omega$-তে একটি সমাপতিত অপেক্ষক দাও।)
\end{prob}

\begin{prob}
হ্রাসকরণ যুক্তির সাহায্যে দেখাও যে স্বাভাবিক সংখ্যার সব ক্রমযুগলের \emph{সেটের} সেট,
অর্থাৎ $\Pow{\Nat \times \Nat}$, !!{nonenumerable}।
\end{prob}

\begin{prob}
হ্রাসকরণ যুক্তির সাহায্যে দেখাও যে স্বাভাবিক সংখ্যার অসীম অনুক্রমগুলির সেট $\Nat^\omega$ !!{nonenumerable}।
\end{prob}

%\begin{prob}
%ধরি, $P$ হলো $\Nat$ থেকে $\{0\}$ সেটে সকল অপেক্ষকের সেট, এবং $Q$ হলো ধনাত্মক পূর্ণসংখ্যার সেট থেকে $\{0\}$ সেটে সকল \emph{আংশিক}
%অপেক্ষকের সেট। দেখাও যে $P$~!!{enumerable}, কিন্তু $Q$ নয়। (ইঙ্গিত: $\Bin^\omega$-কে তালিকায়ন করার সমস্যাটিকে~$Q$-কে তালিকায়ন করার সমস্যায় হ্রাস করো।)
%\end{prob}

\begin{prob}
ধরি, $S$ হলো $\Nat$ থেকে $\{0,1\}$ সেটে সকল !!{surjection}s-এর সেট, অর্থাৎ $S$-এর সদস্য হল সকল !!{surjection}s~$f \colon \Nat
\to \Bin$। দেখাও যে $S$ !!{nonenumerable}।
\end{prob}

\begin{prob}
দেখাও যে সকল বাস্তব সংখ্যার সেট~$\Real$ !!{nonenumerable}।
\end{prob}

\end{document}
